\documentclass[11pt,a4paper]{article}
\usepackage[margin=2.2cm,headheight=16pt]{geometry}
\usepackage[spanish,es-nodecimaldot]{babel}
\usepackage{microtype}
\usepackage{xcolor}
\usepackage{booktabs}
\usepackage{longtable}
\usepackage{tabularx}
\usepackage{array}
\usepackage{enumitem}
\usepackage{fancyhdr}
\usepackage{lastpage}
\usepackage{titlesec}
\usepackage{listings}
\usepackage{amsmath}
\usepackage[most]{tcolorbox}
\usepackage{hyperref}

\definecolor{MemBlue}{HTML}{1F49C4}
\definecolor{MemDark}{HTML}{161A21}
\definecolor{MemGray}{HTML}{4B5563}
\definecolor{MemLight}{HTML}{E7ECFA}
\definecolor{MemGreen}{HTML}{0D7355}
\definecolor{MemOrange}{HTML}{A06B06}
\definecolor{MemRed}{HTML}{AE382C}

\hypersetup{colorlinks=true,linkcolor=MemBlue,urlcolor=MemBlue,
  pdftitle={Manual de usuario de MemLab 16 MiB},pdfauthor={Documentación del proyecto MemLab}}
\pagestyle{fancy}
\fancyhf{}
\lhead{\sffamily\small MemLab 16 MiB}
\rhead{\sffamily\small Manual de usuario}
\cfoot{\sffamily\small Página \thepage\ de \pageref{LastPage}}
\titleformat{\section}{\Large\bfseries\sffamily\color{MemBlue}}{\thesection}{0.7em}{}
\titleformat{\subsection}{\large\bfseries\sffamily\color{MemDark}}{\thesubsection}{0.7em}{}
\setlist{nosep,leftmargin=*}
\setlength{\parindent}{0pt}
\setlength{\parskip}{5pt}
\renewcommand{\arraystretch}{1.18}

\lstdefinestyle{memlab}{basicstyle=\ttfamily\footnotesize,
  backgroundcolor=\color{MemLight!45},frame=single,rulecolor=\color{MemBlue!25},
  breaklines=true,showstringspaces=false,columns=fullflexible}
\lstset{style=memlab}
\newtcolorbox{paso}[1]{colback=MemLight!50,colframe=MemBlue!70!black,
  boxrule=0.7pt,arc=2mm,title={#1},fonttitle=\bfseries\sffamily}
\newtcolorbox{aviso}[1][Importante]{colback=MemOrange!7,colframe=MemOrange!75!black,
  boxrule=0.6pt,arc=2mm,title={#1},fonttitle=\bfseries\sffamily}
\newtcolorbox{resultado}[1][Qué observar]{colback=MemGreen!6,colframe=MemGreen!70!black,
  boxrule=0.6pt,arc=2mm,title={#1},fonttitle=\bfseries\sffamily}
\newcommand{\control}[1]{\textsf{\textbf{#1}}}
\newcommand{\archivo}[1]{\texttt{#1}}

\begin{document}

\begin{titlepage}
  \centering
  \vspace*{2.2cm}
  {\sffamily\bfseries\fontsize{34}{39}\selectfont\color{MemBlue} MemLab 16 MiB\par}
  \vspace{0.5cm}
  {\sffamily\fontsize{23}{28}\selectfont Manual de usuario\par}
  \vspace{0.8cm}
  {\Large Aprende a asignar, liberar y compactar memoria\par}
  \vfill
  \begin{tabular}{rl}
    \textbf{Aplicación:} & simulador web educativo\\
    \textbf{Versión documentada:} & 21 de septiembre de 2026\\
    \textbf{Memoria simulada:} & 16 MiB\\
    \textbf{Instalación:} & no requerida\\
  \end{tabular}
  \vfill
  {\small Este documento describe el comportamiento observado en la versión actual de MemLab.\par}
\end{titlepage}

\tableofcontents
\newpage

\section{¿Qué es MemLab?}

MemLab es un simulador visual para estudiar cómo un sistema operativo reparte la
memoria principal entre varios programas. Trabaja con 16 MiB, desde la dirección
\texttt{0x000000} hasta \texttt{0xFFFFFF}, y permite experimentar sin afectar la
memoria real del equipo.

Con MemLab puedes:

\begin{itemize}
  \item reservar 1, 2 o 4 MiB para el sistema operativo;
  \item comparar particiones fijas, variables y dinámicas;
  \item aplicar primer, mejor y peor ajuste;
  \item cargar o liberar programas y ver dónde quedan ubicados;
  \item estudiar fragmentación interna y externa;
  \item compactar la memoria dinámica;
  \item revisar en la bitácora la explicación de cada decisión.
\end{itemize}

\begin{aviso}
La simulación es educativa. No administra la RAM real, no ejecuta los programas
listados y no guarda cambios al cerrar o recargar la página.
\end{aviso}

\section{Requisitos y apertura}

\subsection{Requisitos}

Solo se necesita un navegador moderno con JavaScript habilitado. No hacen falta
cuenta, conexión a una base de datos, Node.js ni instalación de paquetes. La
conexión a Internet se usa únicamente para descargar tipografías; si no está
disponible, la aplicación sigue funcionando con fuentes locales.

\subsection{Opción rápida: abrir el archivo}

\begin{paso}{Apertura directa}
\begin{enumerate}
  \item Localiza la carpeta del proyecto.
  \item Abre \archivo{index.html} con doble clic.
  \item Espera a que aparezca el mapa de memoria y la tabla de programas.
\end{enumerate}
\end{paso}

\subsection{Opción recomendada: servidor local}

Abre una terminal en la carpeta del proyecto y ejecuta:

\begin{lstlisting}[language=bash]
python3 -m http.server 8000
\end{lstlisting}

Después visita \url{http://localhost:8000/}. Para detener el servidor vuelve a la
terminal y presiona \texttt{Ctrl+C}.

\section{Recorrido por la pantalla}

\subsection{Barra superior}

Muestra el nombre, el rango de direcciones y dos acciones:

\begin{description}
  \item[Escenario de fragmentación.] Prepara automáticamente una distribución con
  procesos y espacios intercalados para estudiar fragmentación.
  \item[Reiniciar.] Recrea los programas y la memoria usando la configuración
  actualmente seleccionada. No necesariamente vuelve al método fijo original.
\end{description}

\subsection{Pestañas de método}

Las tres pestañas superiores cambian el modo de administración. Al cambiar de
método, la memoria se reconstruye y todos los procesos vuelven a la cola.

\begin{longtable}{p{0.25\textwidth}p{0.60\textwidth}}
\toprule
\textbf{Método} & \textbf{Idea principal}\\
\midrule
\endhead
Estáticas de tamaño fijo & Divide la memoria de usuario en partes iguales. Es fácil de entender, pero desperdicia espacio dentro de cada partición ocupada.\\
Estáticas de tamaño variable & Usa una tabla de tamaños diferentes. Permite comparar algoritmos de ajuste, pero las particiones siguen sin poder combinarse.\\
Particiones dinámicas & Crea un bloque del tamaño exacto del programa. Evita fragmentación interna, pero puede dejar huecos externos dispersos.\\
\bottomrule
\end{longtable}

\subsection{Indicadores KPI}

\begin{description}
  \item[Utilización.] Porcentaje de la memoria de usuario ocupado por datos útiles
  de procesos; no incluye el S.O. ni la fragmentación interna.
  \item[Fragmentación interna.] Espacio sobrante dentro de particiones asignadas.
  En modo dinámico siempre es cero.
  \item[Fragmentación externa.] En dinámica, memoria libre fuera del hueco más
  grande; en estáticas, suma de particiones libres que no pueden unirse.
  \item[Mayor bloque libre.] Tamaño máximo de un programa que podría entrar de
  forma inmediata, suponiendo que exista un bloque de ese tamaño.
  \item[Procesos.] Cantidad cargada frente al total; el subtítulo agrupa trabajos
  en cola y rechazados.
\end{description}

\subsection{Mapa de memoria}

La columna representa los 16 MiB completos. La altura de cada bloque es
proporcional a su tamaño. Los patrones distinguen:

\begin{itemize}
  \item sistema operativo;
  \item bloques ocupados, con un color por proceso;
  \item particiones o huecos libres;
  \item fragmentación interna rayada;
  \item zona no particionada, que no puede utilizarse.
\end{itemize}

Haz clic en cualquier bloque para seleccionarlo. El panel \control{Bloque
seleccionado} mostrará su rango, tamaño y, si contiene un proceso, el reparto de
código, datos, pila y montículo.

\subsection{Tabla de programas}

Cada fila contiene nombre, PID, barra de segmentos, tamaño total, dirección base,
estado y acción disponible. Los estados significan:

\begin{description}
  \item[cola:] el trabajo espera una asignación;
  \item[cargado:] tiene memoria y puede liberarse;
  \item[rechazado:] el último intento no encontró un bloque adecuado.
\end{description}

\subsection{Bitácora}

La bitácora coloca el evento más reciente arriba. Explica cambios de configuración,
asignaciones, liberaciones, rechazos, compactaciones e inspecciones. Las marcas
\texttt{t001}, \texttt{t002}, etc. indican el orden lógico, no segundos reales.

\section{Primer uso guiado}

Al abrir la aplicación se seleccionan 2 MiB para el S.O. y siete particiones fijas
de 2 MiB. Se cargan Editor de texto, Compilador C, Terminal y Hoja de cálculo.

\begin{paso}{1. Examinar el estado inicial}
\begin{enumerate}
  \item Observa los cinco KPI.
  \item Haz clic en el bloque del Editor de texto.
  \item Comprueba el rango, el registro base y sus cuatro segmentos.
  \item Identifica la franja de fragmentación interna dentro de su partición.
\end{enumerate}
\end{paso}

\begin{resultado}
El Editor ocupa 1 MiB dentro de una partición de 2 MiB; queda 1 MiB de
fragmentación interna que ningún otro proceso puede aprovechar.
\end{resultado}

\begin{paso}{2. Cargar y liberar}
\begin{enumerate}
  \item Busca un programa en estado \control{cola}.
  \item Presiona \control{Cargar}.
  \item Revisa su base y la nueva línea de bitácora.
  \item Presiona \control{Liberar} y observa cómo vuelve a cola.
\end{enumerate}
\end{paso}

\begin{paso}{3. Reiniciar el ejercicio}
Presiona \control{Reiniciar}. El catálogo y las asignaciones iniciales se recrean
con el método, el S.O. y las opciones que estén activos en ese momento.
\end{paso}

\section{Uso de cada método}

\subsection{Particiones fijas iguales}

\begin{enumerate}
  \item Selecciona \control{Estáticas de tamaño fijo}.
  \item Elige 1, 2 o 4 MiB para el S.O.
  \item Elige entre 2 y 16 particiones según las opciones disponibles.
  \item Lee el tamaño calculado debajo del selector.
  \item Carga programas uno por uno o usa \control{Cargar la cola completa}.
\end{enumerate}

El algoritmo de ajuste no se puede elegir en este modo: como todas las particiones
miden lo mismo, se toma la primera libre capaz de contener el proceso. Si un
programa supera el tamaño de una partición, no podrá cargarse aunque la suma de
varias particiones libres sea suficiente.

\subsection{Particiones fijas variables}

\begin{enumerate}
  \item Selecciona \control{Estáticas de tamaño variable}.
  \item Revisa los chips que indican tamaños en KiB.
  \item Para eliminar una partición, presiona la \(\times\) de su chip.
  \item Para añadir una, escribe al menos 64 KiB y presiona \control{Añadir}.
  \item Escoge primer, mejor o peor ajuste.
  \item Carga trabajos y compara las decisiones en la bitácora.
\end{enumerate}

Los tamaños añadidos se ordenan automáticamente de menor a mayor. La interfaz no
permite añadir una partición si la suma excede la memoria de usuario disponible.
Sin embargo, si primero defines una tabla grande y después aumentas el S.O., algunas
entradas finales pueden no construirse; revisa siempre el mapa y el texto de suma.

\subsection{Particiones dinámicas}

\begin{enumerate}
  \item Selecciona \control{Particiones dinámicas}.
  \item Escoge un algoritmo de ajuste.
  \item Carga varios programas.
  \item Libera procesos intermedios para producir huecos separados.
  \item Observa \control{Frag. externa} y \control{Mayor bloque libre}.
  \item Presiona \control{Compactar ahora} y compara el mapa antes/después.
\end{enumerate}

Activa \control{Compactar automáticamente al fallar una asignación} si quieres que
el simulador compacte cuando el total libre alcanza, pero ningún hueco individual
es suficientemente grande. No compacta si falta memoria total.

\section{Algoritmos de ajuste en palabras simples}

\begin{longtable}{p{0.19\textwidth}p{0.43\textwidth}p{0.27\textwidth}}
\toprule
\textbf{Algoritmo} & \textbf{Cómo decide} & \textbf{Qué conviene observar}\\
\midrule
\endhead
Primer ajuste & Recorre los bloques en orden de dirección y elige el primero que alcance. & Rapidez conceptual y dependencia del orden.\\
Mejor ajuste & Compara todos los candidatos y usa el más pequeño posible. & Resto inmediato pequeño y aparición de huecos diminutos.\\
Peor ajuste & Compara todos y usa el bloque más grande. & Resto grande potencialmente reutilizable.\\
\bottomrule
\end{longtable}

No existe un algoritmo universalmente mejor: el resultado depende del orden y
tamaño de los trabajos. MemLab permite repetir el mismo escenario con distintas
políticas para comparar.

\section{Crear programas personalizados}

Abre el desplegable \control{Añadir un programa a la cola}. Completa:

\begin{itemize}
  \item nombre, con máximo visual de 28 caracteres en el formulario;
  \item código, datos, pila y heap en KiB;
  \item como mínimo 8 KiB de código; los otros segmentos pueden quedar en cero si
  el navegador acepta el valor introducido.
\end{itemize}

Presiona \control{Añadir a la cola}. El tamaño final es la suma de los cuatro
segmentos. No se acepta un proceso mayor que la memoria de usuario actual.

\begin{aviso}[Datos de práctica]
Usa nombres de texto normal. Esta versión representa parte de la bitácora como HTML;
no introduzcas etiquetas ni código en el nombre del programa.
\end{aviso}

\section{Escenario de fragmentación}

El botón \control{Escenario de fragmentación} realiza automáticamente lo siguiente:

\begin{enumerate}
  \item restaura los ocho programas de ejemplo;
  \item reconstruye la memoria con la configuración activa;
  \item intenta cargarlos en orden;
  \item libera los procesos cargados en posiciones alternas;
  \item muestra cuánto queda libre, cuántos bloques hay y cuál es el mayor;
  \item sugiere intentar el trabajo pendiente más grande.
\end{enumerate}

En modo dinámico, este escenario es la forma más rápida de estudiar el contraste
entre memoria libre total y memoria contigua. Prueba a cargar el proceso sugerido,
compacta y vuelve a intentarlo.

\section{Ejercicios recomendados}

\subsection{Ejercicio A: fragmentación interna}

\begin{enumerate}
  \item Usa 2 MiB de S.O. y siete particiones fijas.
  \item Reinicia.
  \item Inspecciona Editor, Terminal y Hoja de cálculo.
  \item Suma las franjas de fragmentación interna.
\end{enumerate}

Resultado esperado inicial: 3\,136 KiB de fragmentación interna entre los cuatro
procesos precargados.

\subsection{Ejercicio B: comparar ajustes}

\begin{enumerate}
  \item Selecciona particiones variables y reinicia.
  \item Anota dónde se ubica cada trabajo con primer ajuste.
  \item Repite desde el mismo estado con mejor ajuste y luego con peor ajuste.
  \item Compara mayor bloque libre, rechazos y fragmentación interna.
\end{enumerate}

\subsection{Ejercicio C: compactación}

\begin{enumerate}
  \item Selecciona dinámica y desactiva compactación automática.
  \item Ejecuta el escenario de fragmentación.
  \item Anota libre total, mayor hueco y fragmentación externa.
  \item Compacta manualmente.
  \item Comprueba que los procesos quedan juntos al inicio y un solo hueco al final.
\end{enumerate}

\section{Programas incluidos}

\begin{longtable}{p{0.40\textwidth}r p{0.39\textwidth}}
\toprule
\textbf{Programa} & \textbf{Tamaño} & \textbf{Comentario práctico}\\
\midrule
\endhead
Editor de texto & 1\,024 KiB & Encaja exactamente en 1 MiB y deja 1 MiB en una partición de 2 MiB.\\
Compilador C & 2\,048 KiB & Encaja exactamente en la partición fija predeterminada.\\
Navegador web & 3\,840 KiB & No cabe en una partición fija predeterminada.\\
Motor de base de datos & 3\,072 KiB & Útil para tablas variables y huecos medianos.\\
Terminal & 448 KiB & Produce fragmentación interna alta en bloques grandes.\\
Render 3D & 6\,144 KiB & Trabajo más grande; exige una partición o hueco de 6 MiB.\\
Servicio de red & 768 KiB & Permite rellenar espacios pequeños.\\
Hoja de cálculo & 1\,536 KiB & Deja 512 KiB en una partición de 2 MiB.\\
\bottomrule
\end{longtable}

\section{Solución de problemas}

\begin{longtable}{p{0.32\textwidth}p{0.55\textwidth}}
\toprule
\textbf{Situación} & \textbf{Qué hacer}\\
\midrule
\endhead
La página aparece sin datos & Confirma que JavaScript esté habilitado y que \archivo{js/script.js} conserve su ruta relativa. Abre la consola del navegador para buscar errores.\\
No se ven los estilos & Verifica que exista \archivo{css/style.css}. Si usas servidor, confirma que responda con estado 200.\\
Un programa es rechazado & Lee la bitácora. Puede ser mayor que cada partición, faltar memoria total o existir fragmentación externa.\\
Hay memoria libre pero no carga & Revisa \control{Mayor bloque libre}. En dinámica, compacta; en estáticas, los bloques no se pueden unir.\\
Compactar no cambia nada & Solo aplica a dinámica. También puede no haber procesos o ya estar contiguos.\\
Cambiar método borró asignaciones & Es el comportamiento esperado: cada disposición tiene límites diferentes y los procesos vuelven a cola.\\
Reiniciar no volvió a siete particiones & Reiniciar conserva la configuración activa; selecciona primero el método y valores deseados.\\
Las letras se ven diferentes & Las fuentes web no cargaron; la aplicación usa una fuente alternativa y sigue funcionando.\\
\bottomrule
\end{longtable}

\section{Accesibilidad y uso con teclado}

Los botones y campos se pueden recorrer con \texttt{Tab}; usa \texttt{Enter} o
barra espaciadora para activarlos. El foco visible indica el control actual. Las
pestañas y algoritmos anuncian su estado mediante \texttt{aria-pressed}, y la
bitácora informa cambios con una región activa cortés. El diseño sigue el tema
claro u oscuro del sistema y evita transiciones si está habilitada la reducción de
movimiento.

En pantallas estrechas, el mapa pasa encima de los paneles laterales. Las tablas
pueden desplazarse horizontalmente.

\section{Glosario}

\begin{description}
  \item[Base.] Primera dirección física asignada a un proceso o bloque.
  \item[Compactación.] Reubicación de procesos para reunir el espacio libre.
  \item[Fragmentación externa.] Espacio libre dividido en regiones separadas.
  \item[Fragmentación interna.] Espacio desperdiciado dentro de una partición ocupada.
  \item[Heap o montículo.] Segmento usado para memoria dinámica del programa.
  \item[Hueco.] Región libre en el método dinámico.
  \item[KiB.] 1\,024 bytes.
  \item[MiB.] 1\,024 KiB o 1\,048\,576 bytes.
  \item[PID.] Identificador numérico interno del proceso simulado.
  \item[Pila.] Segmento asociado a llamadas, variables locales y control de ejecución.
  \item[Registro límite.] Tamaño máximo válido relativo a la base del proceso.
\end{description}

\section{Cierre de la sesión}

No existe una operación de guardado. Si necesitas conservar resultados, toma una
captura o registra manualmente KPI y bitácora antes de cerrar. Detén el servidor
local con \texttt{Ctrl+C} si lo iniciaste desde la terminal.

\end{document}
